\chapter{Model: finite-horizon MDPs, entropic values and best-policy identification}
\label{chap:mdp}

This chapter fixes the objects of the paper (Section~2 of \cite{essakine2026tight}, together
with the exponential Bellman equation of its Section~4) in the form in which they are
formalized on top of the LML framework of sequential algorithms and environments
(\texttt{Learning.Algorithm}, \texttt{Learning.Environment}, \texttt{Learning.IsAlgEnvSeq},
\texttt{Learning.stationaryEnv}) and of LML's MDPs (\texttt{Learning.MDP}). Its definitions are
library material (namespace \texttt{Learning.MDP}, files \texttt{Episodic.lean},
\texttt{Entropic.lean} of \texttt{Essakine2026Tight/LeanMachineLearning/ReinforcementLearning/MDP/}).

\textbf{Indexing convention.} The paper's steps are $h = 1, \dots, H$ with the terminal step
$H + 1$: a state-action pair is visited at the steps $1, \dots, H$ and the state $S_{H+1}$ is
the last state of an episode. The blueprint keeps this $1$-indexing throughout. Lean is
$0$-indexed: the steps of the trajectory laws and of the values are $h \in \N$ (the horizon
$H$ is an argument of the values, and the values are trivial from the step $H$ on), while the
steps of the $H$-step policies played by the learner and of the quantities computed by the
algorithm are $h \in \texttt{Fin H}$ or $h \in \texttt{Fin (H + 1)}$; the paper's step $h$ is
Lean's $h - 1$. The quantity
$e^{\beta(H - h)}$ of the paper (the range of $Z^*_{h+1}$, see Lemma~\ref{lem:opt_exp_value_range})
is $e^{\beta k}$ with $k = H - 1 - h$ the number of steps \emph{after} the Lean step $h$, and the
clip $e^{\beta(H + 1 - h)}$ of Chapter~\ref{chap:algorithm} is $e^{\beta(k + 1)}$. Backward
recursions are indexed in Lean by the number $j$ of steps to the end: $j = 0$ is the terminal
step $H + 1$ and the paper's step $h$ has $j = H + 1 - h$. Episodes are indexed by
$t = 0, 1, \dots$: $t$ is the number of episodes already played and $\pi^{t+1}$ is the policy
computed from them (LML rounds are $0$-indexed: the policy $\pi^{t+1}$ is the action of the
round $t$).

\section{MDPs and policies}
\label{sec:mdp_mdps}

\begin{definition}[Finite MDP]\label{def:mdp}
\lean{MDP, Learning.MDP.env, Learning.MDP.policyAlg, Learning.MDP.policyMeasure}\leanok
A (stationary) \emph{Markov decision process} on a finite state space $\Ss$ and a finite action
space $\As$ is given by a transition kernel $P(\cdot \mid s, a)$, a probability measure on
$\Ss$ for every $(s, a) \in \Ss \times \As$, and a reward kernel $R(\cdot \mid s, a)$, a
probability measure on $\R$ for every $(s, a)$ (the law of the reward received when playing $a$
in $s$). Its \emph{environment} from an initial law $\mu_0$ on $\Ss$ is the LML environment
whose observation at each round is the current state (drawn from $\mu_0$ at the first round
and from $P(\cdot \mid s, a)$ after the round $(s, a)$) and whose feedback to the action $a$ in
the state $s$ is a reward drawn from $R(\cdot \mid s, a)$.
\end{definition}

\textbf{Lean remark.} This is LML's \texttt{Learning.MDP 𝓢 𝓐 𝓡} (fields \texttt{M.P},
\texttt{M.R}, Markov kernels on \texttt{𝓢 × 𝓐}), copied from the \texttt{mdp} branch of LML to
\texttt{Essakine2026Tight/LeanMachineLearning/ReinforcementLearning/MDP/Basic.lean}, and its
environment \texttt{MDP.env M μ₀ : Environment 𝓢 𝓐 𝓡}. The state and action spaces of this
blueprint are finite types with the discrete $\sigma$-algebra (\texttt{Fintype},
\texttt{MeasurableSingletonClass}), so every function on them is measurable and the kernels are
determined by their transition probabilities.

\begin{definition}[Finite-horizon MDP]\label{def:episodic_mdp}
\lean{Learning.MDP.Episodic.EpisodicMDP, Learning.MDP.Episodic.EpisodicMDP.ofDet, Learning.MDP.Episodic.EpisodicMDP.transVec, Learning.MDP.Episodic.EpisodicMDP.meanReward, Learning.MDP.Episodic.EpisodicMDP.RewardsIn}\leanok
\uses{def:mdp}
A \emph{finite-horizon MDP} with horizon $H \ge 1$ is
$\mathcal M = (\Ss, \As, H, (p_h)_{h \le H}, (R_h)_{h \le H})$ where $\Ss$, $\As$ are finite,
$p_h(\cdot \mid s, a)$ is a probability vector on $\Ss$ (the transition kernel at step $h$) and
$R_h(\cdot \mid s, a)$ a probability measure on $\R$ (the reward kernel at step $h$), for
$h \in \{1, \dots, H\}$ and $(s, a) \in \Ss \times \As$. We write $S := \abs{\Ss}$, $A := \abs{\As}$,
$(p_h f)(s, a) := \sum_{s'} p_h(s' \mid s, a) f(s') = \E_{S' \sim p_h(\cdot \mid s, a)}[f(S')]$
for $f : \Ss \to \R$, and $\Var_{p_h}(f)(s, a) := p_h\big((f - (p_h f)(s, a))^2\big)(s, a)$;
more generally $q f := \sum_{s'} q(s') f(s')$ and $\Var_q(f) := q(f - qf)^2$ for a probability
vector $q$ on $\Ss$. The \emph{mean reward} is $r_h(s, a) := \int x\, R_h(dx \mid s, a)$.
\end{definition}

\textbf{Lean remark.} \texttt{EpisodicMDP S A} with fields
\texttt{trans : ℕ → Kernel (S × A) S} and \texttt{reward : ℕ → Kernel (S × A) ℝ}
(Markov kernels indexed by all the steps $h \in \N$, the horizon $H$ being an argument of the
values and of the episode environment; the state and action spaces are arbitrary measurable
spaces, the finiteness of the paper entering through the hypotheses of the statements that
need it; the reward and the next state are independent given $(s, a)$, which is the
paper's model and which the deterministic rewards of Definition~\ref{def:reward_fn} make
harmless). \texttt{M.transVec h s a : S → ℝ} is the vector of transition probabilities
$p_h(\cdot \mid s, a)$ and \texttt{M.meanReward h s a} the mean reward; the operators
$q \mapsto qf$ and $\Var_q$ on finite-state vectors are \texttt{vecExp} and \texttt{vecVar}
(\texttt{Vec.lean}), with \texttt{vecExp (M.transVec h s a) f} for $(p_h f)(s, a)$. A
finite-horizon MDP with the transition kernels \texttt{trans} and a deterministic reward function
\texttt{r} is \texttt{EpisodicMDP.ofDet trans r}.

\begin{definition}[Reward function; the class $\mathcal M_r$]\label{def:reward_fn}
\lean{Learning.MDP.Episodic.EpisodicMDP.HasRewardFn}\leanok
\uses{def:episodic_mdp}
A \emph{reward function} is a map $r : \{1, \dots, H\} \times \Ss \times \As \to [0, 1]$,
written $r_h(s, a)$. The MDP $\mathcal M$ \emph{has reward function $r$} if all its rewards are
deterministic with values given by $r$: $R_h(\cdot \mid s, a) = \delta_{r_h(s, a)}$ for all $h$,
$s$, $a$. For a fixed reward function $r$, $\mathcal M_r$ denotes the class of all finite-horizon
MDPs on $(\Ss, \As, H)$ with reward function $r$ (their transition kernels are arbitrary). All
the MDPs of the paper are in some $\mathcal M_r$: the rewards are known, deterministic and in
$[0, 1]$; only the transitions are unknown to the learner.
\end{definition}

\textbf{Lean remark.} \texttt{M.HasRewardFn r} is the predicate
$\forall h\, s\, a,\ \texttt{M.reward h (s, a) = Measure.dirac (r h s a)}$; the class
$\mathcal M_r$ is the subtype \texttt{\{M : EpisodicMDP S A // M.HasRewardFn r\}}, with
\texttt{r : ℕ → S → A → ℝ}. The
condition $r \in [0, 1]$ is a separate hypothesis \texttt{∀ h s a, r h s a ∈ Set.Icc 0 1} of
the statements that need it (all the results of Chapters~\ref{chap:algorithm}
to~\ref{chap:lower_bound}).

\begin{definition}[Deterministic policies]\label{def:policy}
\lean{Learning.MDP.Episodic.Policy}\leanok
\uses{def:episodic_mdp}
A \emph{(deterministic, non-stationary) policy} is $\pi = (\pi_h)_{h \le H}$ with
$\pi_h : \Ss \to \As$; $\pi_h(s)$ is the action played in the state $s$ at step $h$. There are
finitely many policies ($A^{SH}$). For $g : \Ss \times \As \to \R$ we write
$(\pi_h g)(s) := g(s, \pi_h(s))$ for the composition with the policy at step $h$.
\end{definition}

\textbf{Lean remark.} The trajectory laws and the values are defined for a policy
\texttt{π : ℕ → S → A} on all the steps. As for a stochastic process in Mathlib, the policy is not
bundled with its measurability: the definitions are total (the trajectory law of a policy which
is not measurable at every step is the zero measure, as \texttt{Measure.map} of a non-measurable
function), and the lemmas take the hypothesis \texttt{hπ : ∀ h, Measurable (π h)}. The learner
plays $H$-step policies \texttt{Policy S A H := Fin H → S → A}, a finite type (\texttt{Fintype}
instance by \texttt{Pi.instFintype}) with the discrete $\sigma$-algebra, extended to all the steps
by \texttt{Policy.extend} (the policy on the first $H$ steps, a fixed action beyond;
\texttt{Policy.extend\_fin}, \texttt{Policy.extend\_injective}, and
\texttt{Policy.measurable\_extend} on a countable discrete state space); the values of
\texttt{π.extend} only depend on \texttt{π} (\texttt{expValue\_congr}). The paper also allows
stochastic policies in the notation $(\pi_h g)(s) = \E_{A \sim \pi_h(\cdot \mid s)} g(s, A)$;
only deterministic policies are needed (the learner plays deterministic policies and the
optimal policies are deterministic, Lemma~\ref{lem:opt_bellman}).

\section{Trajectory laws through the time-augmented MDP}
\label{sec:mdp_trajectories}

The trajectory of a policy in a finite-horizon MDP is defined as the trajectory of a
stationary policy in a stationary MDP, the \emph{time-augmented} MDP, whose state records the
current step; this reuses LML's trajectory measures and makes the Markov property a consequence
of the Ionescu--Tulcea construction.

\begin{definition}[Time-augmented MDP]\label{def:layered_mdp}
\lean{Learning.MDP.Episodic.EpisodicMDP.augmented, Learning.MDP.Episodic.EpisodicMDP.augTrans, Learning.MDP.Episodic.EpisodicMDP.augReward, Learning.MDP.Episodic.augPolicy, Learning.MDP.Episodic.measurable_augPolicy, Learning.MDP.Episodic.Policy.extend, Learning.MDP.Episodic.Policy.measurable_extend}\leanok
\uses{def:episodic_mdp, def:policy, def:mdp}
The \emph{time-augmented MDP} of $\mathcal M$ is the stationary MDP (Definition~\ref{def:mdp})
on the state space $\Ss \times \N$ with actions $\As$: from a state $(s, h)$ and an action $a$,
the reward has law $R_h(\cdot \mid s, a)$ and the next state is $(s', h + 1)$ with
$s' \sim p_h(\cdot \mid s, a)$. A policy $\pi = (\pi_h)_{h \in \N}$ of $\mathcal M$ becomes the
stationary policy $\pi^{\mathrm{aug}}(s, h) := \pi_h(s)$; an $H$-step policy is extended by a
fixed action beyond the horizon.
\end{definition}

\textbf{Lean remark.} \texttt{M.augmented : MDP (S × ℕ) A ℝ} (kernels \texttt{M.augTrans},
\texttt{M.augReward}) and \texttt{augPolicy π : S × ℕ → A}, measurable when every
\texttt{π h} is (\texttt{measurable\_augPolicy}, by
\texttt{measurable\_from\_prod\_countable\_left}). The steps are all of $\N$: there is no
terminal layer, the trajectory runs forever and the finite horizon only enters through the
return of the first $H - h$ rounds (Definition~\ref{def:episode_return}). An $H$-step policy
\texttt{π : Policy S A H} is \texttt{π.extend}, which plays a default action of the nonempty
type \texttt{A} beyond the horizon.

\begin{definition}[Trajectory law from a state at a step]\label{def:step_law}
\lean{Learning.MDP.Episodic.stepLaw}\leanok
\uses{def:layered_mdp}
For a policy $\pi$, a state $s$ and a step $h \in \N$, the \emph{trajectory law}
$\Prob^\pi_{(s, h)}$ is the law of the trajectory of the stationary policy $\pi^{\mathrm{aug}}$ in
the time-augmented MDP started at the state $(s, h)$, i.e.\ the sequence of rounds
$((S_{h+j}, h + j), A_{h+j}, R_{h+j})_{j \ge 0}$ (observation, action, feedback). We write
$(S_i, A_i, R_i)_{i \ge h}$ for this trajectory (so $S_h = s$, $A_i = \pi_i(S_i)$ and
$R_i \sim R_i(\cdot \mid S_i, A_i)$ for $i \ge h$) and $\E^\pi_{(s, h)}$ for the expectation under
$\Prob^\pi_{(s, h)}$; $\Prob^\pi := \Prob^\pi_{(s_1, 1)}$ and $\E^\pi := \E^\pi_{(s_1, 1)}$ when an
initial state $s_1$ is fixed.
\end{definition}

\textbf{Lean remark.} \texttt{stepLaw M π (s, h) := M.augmented.policyMeasure (augPolicy π) \_ (s, h)}
for \texttt{π : ℕ → S → A} measurable at every step (\texttt{stepLaw\_of\_measurable}; it is
\texttt{0} otherwise) and \texttt{h : ℕ}, a probability measure
(\texttt{isProbabilityMeasure\_stepLaw}) on
\texttt{ℕ → Round (S × ℕ) A ℝ} (LML \texttt{MDP.policyMeasure}, the trajectory measure
\texttt{trajMeasure} of the stationary policy algorithm \texttt{policyAlg} in
\texttt{MDP.env M.augmented (Measure.dirac (s, h))}); \texttt{stepLawKernel M π h} is the same
law as a kernel from the state. The round $j$ of the trajectory corresponds to the paper's step
$h + j$.

\begin{definition}[Episode return and state sequence]\label{def:episode_return}
\lean{Learning.MDP.Episodic.episodeReturn, Learning.MDP.Episodic.episodeStates, Learning.MDP.Episodic.Traj}\leanok
\uses{def:step_law}
The \emph{return from step $h$} of a trajectory is $R^\pi_h := \sum_{i = h}^H R_i$, the sum of
the first $H + 1 - h$ rewards of the trajectory from $(s, h)$ ($0$ if $h > H$). The
\emph{state sequence} of an episode started at $(s_1, 1)$ is $(S_1, \dots, S_{H+1})$: the
observation states of the first $H + 1$ rounds. We call the pair (state sequence, rewards) the
\emph{episode}.
\end{definition}

\textbf{Lean remark.} \texttt{episodeReturn n : (ℕ → Round (S × ℕ) A ℝ) → ℝ},
\texttt{traj ↦ ∑ k ∈ range n, (traj k).feedback}, the sum of the first \texttt{n} rewards: the
return from the ($0$-indexed) step \texttt{h} is \texttt{episodeReturn (H - h)}, and
\texttt{episodeStates H : \_ → Traj S H}, \texttt{traj ↦ fun i : Fin (H + 1) ↦ (traj i).obs.1},
where \texttt{Traj S H := Fin (H + 1) → S}. Both are measurable (finite sums and coordinate
projections).

\begin{lemma}[Markov property of the trajectory of a stationary policy]\label{lem:policy_markov}
\lean{Learning.MDP.hasCondDistrib_shiftRound_trajMeasure, Learning.MDP.hasLaw_step_zero_policyMeasure, Learning.MDP.hasCondDistrib_obs_one_trajMeasure}\leanok
\uses{def:mdp}
Let $M$ be an MDP whose state space has measurable singletons, $\pi$ a measurable
stationary deterministic policy and $\mu_0$ an initial law, and let $(O_j, A_j, Y_j)_{j \ge 0}$ be
the canonical trajectory of $\pi$ in the environment of $M$ from $\mu_0$. Then the first round
$(O_0, A_0, Y_0)$ has law $\mu_0(ds)\,\delta_{\pi(s)}(da)\,R(dy \mid s, a)$, the state $O_1$ has
conditional law $P(\cdot \mid O_0, A_0)$ given the first round, and the shifted trajectory
$(O_{j+1}, A_{j+1}, Y_{j+1})_{j \ge 0}$ has conditional law $\Prob^\pi_{O_1}$ (the trajectory law of
$\pi$ from the state $O_1$) given $((O_0, A_0, Y_0), O_1)$.
\end{lemma}

\begin{proof}
\leanok
Let $Q$ be the law of $((O_0, A_0, Y_0), O_1)$ followed by an independent trajectory drawn from
$\Prob^\pi_{O_1}$, a measure on (round, state, trajectory), and let $\Phi$ map such a triple to the
trajectory obtained by prepending the round to the trajectory. Under $Q$, the process
$\Phi$ has the conditional laws of the canonical trajectory: its round $0$ has the law of the
first round; its round $1$, given round $0$, has observation $P(\cdot \mid O_0, A_0)$ and then the
action and reward of $\pi$ from that state (the initial round of $\Prob^\pi_{O_1}$); its round
$j + 2$, given the previous rounds, has the step kernel of the pair at round $j + 1$, which is the
step kernel of $\Prob^\pi_{O_1}$ at round $j + 1$ because the step kernels at positive rounds do
not depend on the initial law. By uniqueness of the Ionescu--Tulcea measure with given
conditional laws (LML \texttt{Kernel.hasLaw\_trajMeasureFin}), $Q \circ \Phi^{-1}$ is the law of
the canonical trajectory. Since the observation of round $0$ of $\Prob^\pi_{s'}$ is $s'$ almost
surely, $\Phi$ is $Q$-almost surely inverted by (first round, state of round $1$, shift), so the
law of this triple under the trajectory law is $Q$, which is the claim.
\end{proof}

\textbf{Lean remark.} File \texttt{MDP/Markov.lean}, in LML generality (\texttt{Learning.MDP},
arbitrary state space with \texttt{MeasurableSingletonClass}). The trajectory law of $\pi$ as a
kernel from the initial state is \texttt{policyKernel M hπ}, built with Mathlib's
\texttt{Kernel.traj} from the step kernels (\texttt{iicStepKernel}) so that its measurability
is Mathlib's; the first round and the state of round $1$ are
\texttt{firstRoundState}, the shift is \texttt{shiftRound}, the auxiliary measure is
\texttt{shiftMeasure} and the prepending map is \texttt{consRound}. The general
\texttt{HasCondDistrib} tools used are in \texttt{Mathlib/Probability/HasCondDistrib.lean}
(\texttt{hasCondDistrib\_snd\_compProd}, \texttt{HasCondDistrib.comp\_hasLaw},
\texttt{HasCondDistrib.compProd\_left}).

\begin{lemma}[Markov property of the trajectory law]\label{lem:step_law_markov}
\lean{Learning.MDP.Episodic.hasLaw_stepLaw_succ, Learning.MDP.Episodic.ae_obs_zero_stepLaw, Learning.MDP.Episodic.ae_action_zero_stepLaw, Learning.MDP.Episodic.hasLaw_feedback_zero_stepLaw, Learning.MDP.Episodic.ae_feedback_eq_policy_stepLaw, Learning.MDP.Episodic.episodeReturn_succ_eq_add_shiftRound, Learning.MDP.Episodic.ae_stepLaw_succ_of_ae, Learning.MDP.Episodic.lintegral_stepLaw_succ, Learning.MDP.Episodic.integral_stepLaw_succ}\leanok
\uses{def:step_law, def:layered_mdp}
Let $h \in \N$, $s \in \Ss$ and $\pi$ a policy. Under $\Prob^\pi_{(s, h)}$:
\begin{enumerate}
  \item[(i)] the first round is $S_h = s$, $A_h = \pi_h(s)$, and $(R_h, S_{h+1})$ has law
        $R_h(\cdot \mid s, \pi_h(s)) \otimes p_h(\cdot \mid s, \pi_h(s))$;
  \item[(ii)] the conditional law of the shifted trajectory $(S_i, A_i, R_i)_{i \ge h+1}$ given
        the first round is $\Prob^\pi_{(S_{h+1}, h + 1)}$: for every measurable bounded (or
        nonnegative) $F$ on trajectories,
        \[
          \E^\pi_{(s, h)}\big[F\big((S_i, A_i, R_i)_{i \ge h + 1}\big)\big]
          = \E_{(R, s') \sim R_h(\cdot \mid s, \pi_h(s)) \otimes p_h(\cdot \mid s, \pi_h(s))}
            \Big[ \E^\pi_{(s', h + 1)}[F] \Big],
        \]
        and more generally
        $\E^\pi_{(s, h)}[G(R_h, S_{h+1})\, F((S_i, A_i, R_i)_{i \ge h+1})]
        = \E_{(R, s')}[G(R, s')\, \E^\pi_{(s', h+1)}[F]]$ for bounded measurable $G$.
\end{enumerate}
\end{lemma}

\begin{proof}
\uses{def:step_law, def:layered_mdp, lem:policy_markov}
\leanok
The trajectory measure of an algorithm in an environment is the Ionescu--Tulcea measure built
from the initial law of the observation and the step kernels (LML \texttt{trajMeasure}). For the
stationary policy $\pi^{\mathrm{aug}}$ in the time-augmented MDP started at $(s, h)$: the
observation of the round $0$ is $(s, h)$ almost surely (initial law $\delta_{(s, h)}$), the action
is the deterministic $\pi^{\mathrm{aug}}(s, h) = \pi_h(s)$ (the policy algorithm is deterministic),
and the feedback and next observation are drawn from the augmented step kernel at
$((s, h), \pi_h(s))$,
which is $R_h(\cdot \mid s, \pi_h(s))$ for the reward and $s' \mapsto (s', h + 1)$ with
$s' \sim p_h(\cdot \mid s, \pi_h(s))$ for the next state, independently: this is (i). For (ii), the
canonical trajectory process is an algorithm-environment sequence for the policy algorithm and
the augmented environment (LML \texttt{isAlgEnvSeq\_trajMeasure}); the conditional law of the
future rounds $j \ge 1$ given the round $0$ is, by the Ionescu--Tulcea construction, the
trajectory measure of the same algorithm-environment pair started at the observation of round
$1$, that is $(S_{h+1}, h + 1)$ (the policy algorithm and the MDP environment are stationary:
their kernels at round $j$ only depend on the last observation and action, not on $j$ nor on
the earlier rounds). Since the observation of round $1$ is $(S_{h+1}, h + 1)$ and the law of
$(R_h, S_{h+1})$ is the product kernel of (i), the displayed identities follow by integrating
the conditional law, first for $F$ an indicator of a measurable rectangle of the trajectory
space and then by the monotone class theorem.
\end{proof}

\textbf{Lean remark.} File \texttt{MDP/StepLaw.lean}, derived from
Lemma~\ref{lem:policy_markov} for the time-augmented MDP. For a Lean step \texttt{h : ℕ} (the
paper's $h + 1$), \texttt{hasLaw\_stepLaw\_succ} states that under \texttt{stepLaw M π (s, h)}
the triple (reward of round $0$, state of round $1$, shifted trajectory) has law
\texttt{(M.reward h (s, π h s)).prod (M.trans h (s, π h s)) ⊗ₘ Kernel.prodMkLeft ℝ (stepLawKernel M π (h + 1))},
which contains (i) and (ii); the almost sure statements and integrals along one step are
\texttt{ae\_stepLaw\_succ\_of\_ae}, \texttt{lintegral\_stepLaw\_succ} and
\texttt{integral\_stepLaw\_succ}; with a reward function the first reward is
$r_h(s, \pi_h(s))$ almost surely (\texttt{ae\_feedback\_eq\_policy\_stepLaw}). The return of
$n + 1$ rounds is the first reward plus the return of $n$ rounds of the shifted trajectory
(\texttt{episodeReturn\_succ\_eq\_add\_shiftRound}, an identity of functions): this is
$R^\pi_h = R_h + R^\pi_{h+1}$, used by every backward recursion, and the recursions stop at the
step $H$, where the return is the empty sum (no terminal layer is needed).

\begin{definition}[Additive values]\label{def:eval_value}
\uses{def:step_law, def:episode_return}
For $f : \{1, \dots, H\} \times \Ss \times \As \to \R$, a policy $\pi$, a step $h$ and a state $s$,
the \emph{additive value} of $\pi$ for $f$ is
$\E^\pi_{(s, h)}\big[\sum_{i = h}^H f_i(S_i, A_i)\big]$; with $f = r$ (a reward function) it is the
usual (risk-neutral) value $\E^\pi_{(s, h)}[R^\pi_h]$ of the MDP with reward function $r$.
\end{definition}

\textbf{Lean remark.} Not formalized (it would be
\texttt{∫ traj, ∑ k ∈ range (H - h), f (h + k) (traj k).obs.1 (traj k).action ∂(stepLaw M π (s, h))}).
It is not needed for the results of the paper (the entropic values are not additive) and
serves as the risk-neutral reference point.

\section{Exponential and entropic values}
\label{sec:mdp_values}

Throughout, $\beta \ne 0$ is the risk parameter. The paper works in the \emph{exponential space}:
all the recursions are for $Z = e^{\beta V}$, which satisfies a linear Bellman equation, and the
entropic values are recovered by a logarithm.

\begin{definition}[Exponential values]\label{def:exp_value}
\lean{Learning.MDP.Episodic.expValue, Learning.MDP.Episodic.expQ}\leanok
\uses{def:step_law, def:episode_return, def:episodic_mdp}
For $\beta \ne 0$, a policy $\pi$, a step $h \in \{1, \dots, H + 1\}$ and a state $s$, the
\emph{exponential value} is
\[
  Z^\pi_h(s) := \E^\pi_{(s, h)}\big[e^{\beta R^\pi_h}\big] ,
  \qquad\text{so that } Z^\pi_{H+1}(s) = 1 ,
\]
and for $h \le H$ and $a \in \As$ the \emph{exponential state-action value} is
\[
  U^\pi_h(s, a) := \E_{(R, S') \sim R_h(\cdot \mid s, a) \otimes p_h(\cdot \mid s, a)}
    \big[e^{\beta R}\, Z^\pi_{h+1}(S')\big] .
\]
When $\mathcal M$ has the reward function $r$ (Definition~\ref{def:reward_fn}),
$U^\pi_h(s, a) = e^{\beta r_h(s, a)}\,(p_h Z^\pi_{h+1})(s, a)$.
\end{definition}

\textbf{Lean remark.} \texttt{expValue M H β π h s := ∫ traj, Real.exp (β * episodeReturn (H - h) traj) ∂(stepLaw M π (s, h))}
for \texttt{h : ℕ} and \texttt{π : ℕ → S → A}, and
\texttt{expQ M H β π h s a := (∫ x, Real.exp (β * x) ∂M.reward h (s, a)) * ∫ s', expValue M H β π (h + 1) s' ∂M.trans h (s, a)};
with a reward function, \texttt{expQ = Real.exp (β * r h s a) * ∫ s', expValue M H β π (h + 1) s' ∂M.trans h (s, a)}
(\texttt{expQ\_of\_hasRewardFn}), and on a finite state space
\texttt{= Real.exp (β * r h s a) * vecExp (M.transVec h s a) (expValue M H β π (h + 1))}
(\texttt{expQ\_of\_hasRewardFn\_eq\_vecExp}). The terminal value $Z^\pi_{H+1} = 1$ is a lemma
(\texttt{expValue\_of\_le}: the return of $H - h = 0$ rounds is $0$ for $h \ge H$), not a
separate definition.

\begin{definition}[Entropic values]\label{def:entropic_value}
\lean{Learning.MDP.Episodic.entropicValue, Learning.MDP.Episodic.entropicQ}\leanok
\uses{def:exp_value}
The \emph{entropic value} and \emph{entropic state-action value} of $\pi$ are
\[
  V^\pi_h(s) := \frac1\beta \log Z^\pi_h(s)
  = \frac1\beta \log \E^\pi_{(s, h)}\Big[\exp\Big(\beta \sum_{i = h}^H R_i\Big)\Big],
  \qquad
  Q^\pi_h(s, a) := \frac1\beta \log U^\pi_h(s, a) ,
\]
so that $V^\pi_{H+1} = 0$ and $Z^\pi_h = e^{\beta V^\pi_h}$, $U^\pi_h = e^{\beta Q^\pi_h}$ (the
logarithms are well defined by Lemma~\ref{lem:exp_value_range}).
\end{definition}

\textbf{Lean remark.} \texttt{entropicValue M H β π h s := β⁻¹ * Real.log (expValue M H β π h s)},
\texttt{entropicQ}. The entropic risk measure with $\beta > 0$ is risk-seeking and with
$\beta < 0$ risk-averse \cite{borkar2001risk}; the formalization treats both signs with the
same definitions.

\begin{lemma}[Exponential Bellman equation]\label{lem:exp_bellman}
\lean{Learning.MDP.Episodic.expValue_succ, Learning.MDP.Episodic.expValue_of_le, Learning.MDP.Episodic.lexpValue_succ, Learning.MDP.Episodic.lexpValue_lt_top, Learning.MDP.Episodic.expQ_of_hasRewardFn, Learning.MDP.Episodic.expQ_of_hasRewardFn_eq_vecExp, Learning.MDP.Episodic.integrable_exp_mul_of_rewardsIn, Learning.MDP.Episodic.expValue_congr}\leanok
\uses{def:exp_value, lem:step_law_markov}
Assume that the rewards are bounded: $R_h(\cdot \mid s, a)$ is supported in a fixed interval
$[a, b]$ for all $h, s, a$ (true for a reward function with values in $[0, 1]$). Then
$Z^\pi_{h} = 1$ for $h \ge H + 1$ and for every policy $\pi$, step $h \le H$ and state $s$,
\[
  Z^\pi_h(s) = U^\pi_h(s, \pi_h(s)),
\]
i.e.\ $Z^\pi_h(s) = e^{\beta r_h(s, \pi_h(s))}\,(p_h Z^\pi_{h+1})(s, \pi_h(s))$ when $\mathcal M$
has the reward function $r$.
\end{lemma}

\begin{proof}
\uses{lem:step_law_markov, def:exp_value, def:episode_return}
\leanok
Under $\Prob^\pi_{(s, h)}$, $R^\pi_h = R_h + R^\pi_{h+1}$ where $R^\pi_{h+1} = \sum_{i \ge h+1} R_i$
is a function of the shifted trajectory, so
$e^{\beta R^\pi_h} = e^{\beta R_h}\, e^{\beta R^\pi_{h+1}}$. Lemma~\ref{lem:step_law_markov}(ii) with
$G(R, s') = e^{\beta R}$ and $F = e^{\beta R^\pi_{h+1}}$ (bounded, the rewards being bounded)
gives
$Z^\pi_h(s) = \E_{(R, S') \sim R_h \otimes p_h(\cdot \mid s, \pi_h(s))}[e^{\beta R} Z^\pi_{h+1}(S')]
= U^\pi_h(s, \pi_h(s))$, since $\E^\pi_{(s', h+1)}[e^{\beta R^\pi_{h+1}}] = Z^\pi_{h+1}(s')$ by
definition. The deterministic-reward form is Definition~\ref{def:exp_value}.
\end{proof}

\textbf{Lean remark.} The identity is first proved for the Lebesgue integrals
$L^\pi_h(s) := \int^- e^{\beta R^\pi_h}\,d\Prob^\pi_{(s,h)} \in [0, \infty]$ (\texttt{lexpValue}),
without any hypothesis: $L^\pi_h(s) = \big(\int^- e^{\beta x}R_h(dx \mid s, a)\big)\sum_{s'} p_h(s' \mid s, a)L^\pi_{h+1}(s')$
with $a = \pi_h(s)$ (\texttt{lexpValue\_succ}). The boundedness hypothesis
(\texttt{M.RewardsIn (Set.Icc a b)}) makes all $L^\pi_h(s)$ finite
(\texttt{lexpValue\_lt\_top}), and then $Z^\pi_h = (L^\pi_h)^{\mathrm{toReal}}$ satisfies the
Bellman equation (\texttt{expValue\_succ}). Without such a hypothesis the Bochner-integral
identity can fail (a non-integrable exponential return has Bochner integral $0$), so the
hypothesis is not an artifact. The backward inductions on the step \texttt{h : ℕ} are
\texttt{horizon\_induction} (the case $h \ge H$, then the step from $h + 1$ to $h < H$).

\begin{lemma}[Ranges of the exponential values]\label{lem:exp_value_range}
\lean{Learning.MDP.Episodic.expValue_mem_Icc, Learning.MDP.Episodic.expQ_mem_Icc, Learning.MDP.Episodic.ae_episodeReturn_mem_Icc, Learning.MDP.Episodic.expValue_pos}\leanok
\uses{def:exp_value, lem:step_law_markov, lem:exp_bellman}
Assume that the rewards of $\mathcal M$ are in $[0, 1]$ ($R_h(\cdot \mid s, a)$ is supported in
$[0, 1]$ for all $h, s, a$; in particular when $\mathcal M$ has a reward function). Then for every
policy $\pi$, step $h \le H + 1$ and state $s$, $R^\pi_h \in [0, H + 1 - h]$
$\Prob^\pi_{(s, h)}$-a.s., and
\[
  Z^\pi_h(s) \in \big[\min(1, e^{\beta(H + 1 - h)}),\ \max(1, e^{\beta(H + 1 - h)})\big],
  \qquad
  U^\pi_h(s, a) \in \big[\min(1, e^{\beta(H + 1 - h)}),\ \max(1, e^{\beta(H + 1 - h)})\big]
\]
(the second for $h \le H$); in particular $Z^\pi_h(s) > 0$ and $U^\pi_h(s, a) > 0$. For
$\beta > 0$ the intervals are $[1, e^{\beta(H+1-h)}]$, for $\beta < 0$ they are
$[e^{\beta(H+1-h)}, 1]$.
\end{lemma}

\begin{proof}
\uses{lem:step_law_markov, def:exp_value, def:episode_return}
\leanok
By Lemma~\ref{lem:step_law_markov}(i) applied at each round (induction on the rounds, using (ii)
to move to the next step), the reward $R_i$ of the round at step $i \le H$ has law
$R_i(\cdot \mid S_i, A_i)$, hence lies in $[0, 1]$ a.s., and the rewards of the rounds at
steps $i > H$ are $0$. So $R^\pi_h = \sum_{i = h}^H R_i \in [0, H + 1 - h]$ a.s. The function
$x \mapsto e^{\beta x}$ is monotone, so $e^{\beta R^\pi_h}$ lies a.s.\ between $1 = e^{0}$ and
$e^{\beta(H + 1 - h)}$, and so does its expectation $Z^\pi_h(s)$ (monotonicity of the integral of a
probability measure). For $U^\pi_h(s, a)$: $e^{\beta R}$ lies between $1$ and $e^\beta$ for
$R \in [0, 1]$, $Z^\pi_{h+1}(S')$ lies between $1$ and $e^{\beta(H - h)}$ by the first part at the
step $h + 1$, and the product of two numbers lying respectively between $1$ and $e^{\beta a}$
and between $1$ and $e^{\beta b}$ ($a, b \ge 0$) lies between $1$ and $e^{\beta(a + b)}$; the
expectation preserves the bounds. Positivity: the lower bounds are positive.
\end{proof}

\textbf{Lean remark.} The ranges are stated with \texttt{min 1 (Real.exp (β * (H - h)))} and
\texttt{max 1 (Real.exp (β * (H - h)))} for the Lean step \texttt{h : ℕ} (\texttt{H - h} steps to
go, truncated subtraction; for \texttt{expQ} at \texttt{h < H} the same exponent). The
Lean proof of the ranges of $Z^\pi$ and $U^\pi$ is analytic, by backward induction on the
exponential Bellman equation (Lemma~\ref{lem:exp_bellman}), with the product rule
\texttt{mul\_mem\_Icc\_min\_max\_exp}; the almost sure bound on the return is proved by
induction on the number of rounds with \texttt{episodeReturn\_succ\_eq\_add\_shiftRound} and
\texttt{ae\_stepLaw\_succ\_of\_ae}. Positivity \texttt{expValue\_pos} only needs bounded rewards.

\begin{lemma}[Monotonicity of the exponential Bellman operator]\label{lem:exp_value_monotone}
\lean{Learning.MDP.vecExp_mono, Learning.MDP.vecExp_nonneg, Learning.MDP.vecExp_const, Learning.MDP.vecExp_add, Learning.MDP.vecExp_const_mul, Learning.MDP.vecExp_mem_Icc, Learning.MDP.abs_vecExp_le, Learning.MDP.transVec_nonneg, Learning.MDP.sum_transVec, Learning.MDP.Episodic.EpisodicMDP.sum_transVec, Learning.MDP.Episodic.integral_eq_vecExp_transVec}\leanok
\uses{def:episodic_mdp}
For every step $h \le H$ and $(s, a)$: $f \mapsto (p_h f)(s, a)$ is linear on $\R^{\Ss}$,
positive ($f \ge 0 \Rightarrow p_h f \ge 0$), monotone ($f \le g \Rightarrow p_h f \le p_h g$),
$p_h \mathbf 1 = 1$, and $\abs{(p_h f)(s, a)} \le \max_{s'} \abs{f(s')}$. Consequently, for every
$c \in \R$ (in particular $c = e^{\beta r_h(s, a)} > 0$), $f \le g$ implies
$e^{c}(p_h f)(s, a) \le e^{c}(p_h g)(s, a)$; the same holds for every probability vector $q$ on
$\Ss$ in place of $p_h(\cdot \mid s, a)$.
\end{lemma}

\begin{proof}
\uses{def:episodic_mdp}
\leanok
$(p_h f)(s, a) = \sum_{s'} p_h(s' \mid s, a) f(s')$ with nonnegative weights summing to $1$:
linearity, positivity and $p_h \mathbf 1 = 1$ are immediate, monotonicity is positivity applied
to $g - f$, and the bound is the triangle inequality. Multiplying by $e^c > 0$ preserves the
order.
\end{proof}

\textbf{Lean remark.} These are the lemmas of \texttt{Vec.lean}, stated for a vector
\texttt{p : S → ℝ} with the hypotheses \texttt{∀ s, 0 ≤ p s} and \texttt{∑ s, p s = 1} where
needed; \texttt{M.transVec h s a} satisfies both (\texttt{transVec\_nonneg},
\texttt{EpisodicMDP.sum\_transVec}), and \texttt{integral\_eq\_vecExp\_transVec} identifies
$\int f\,dp_h(\cdot \mid s, a)$ with \texttt{vecExp (M.transVec h s a) f}.

\begin{definition}[Optimal values]\label{def:opt_value}
\lean{Learning.MDP.Episodic.optEntropicValue, Learning.MDP.Episodic.optExpValue, Learning.MDP.Episodic.optExpQ}\leanok
\uses{def:entropic_value, def:exp_value}
The \emph{optimal entropic value} is $V^*_h(s) := \sup_\pi V^\pi_h(s)$ (a supremum over all
measurable policies, finite for bounded rewards), the \emph{optimal exponential value} is $Z^*_h(s) := e^{\beta V^*_h(s)}$, and the
\emph{optimal exponential state-action value} is, for $h \le H$,
\[
  U^*_h(s, a) := \E_{(R, S') \sim R_h(\cdot \mid s, a) \otimes p_h(\cdot \mid s, a)}
    \big[e^{\beta R} Z^*_{h+1}(S')\big]
  = e^{\beta r_h(s, a)}\,(p_h Z^*_{h+1})(s, a) \text{ for a reward function } r .
\]
The paper writes $Q^*_h := \beta^{-1} \log U^*_h$; then $Z^*_{H+1} = 1$ and $V^*_{H+1} = 0$.
\end{definition}

\textbf{Lean remark.} \texttt{optEntropicValue M H β h s := ⨆ π : \{π : ℕ → S → A // ∀ k, Measurable (π k)\}, entropicValue M H β π h s}
(a \texttt{ciSup} in $\R$ over the policies measurable at every step, so that the zero measure
of a non-measurable policy never enters the supremum; bounded above for bounded rewards:
\texttt{bddAbove\_entropicValue}; it is attained, Lemma~\ref{lem:opt_bellman}),
\texttt{optExpValue M H β h s := Real.exp (β * optEntropicValue M H β h s)},
\texttt{optExpQ M H β h s a} defined as \texttt{expQ} with \texttt{optExpValue M H β (h + 1)} in
place of the value of a policy.

\begin{lemma}[Optimal exponential value as an extremum]\label{lem:opt_exp_value_eq}
\lean{Learning.MDP.Episodic.expValue_le_optExpValue, Learning.MDP.Episodic.optExpValue_le_expValue, Learning.MDP.Episodic.expValue_eq_exp_mul_entropicValue, Learning.MDP.Episodic.entropicValue_le_optEntropicValue, Learning.MDP.Episodic.bddAbove_entropicValue, Learning.MDP.Episodic.entropicValue_le_of_pos, Learning.MDP.Episodic.entropicValue_le_of_neg}\leanok
\uses{def:opt_value, lem:exp_value_range}
Assume $\beta \ne 0$ and that the rewards are in $[0, 1]$. Then for every $h$ and $s$ the values
$V^\pi_h(s)$ are bounded above uniformly in $\pi$ (by $H - h$ for $\beta > 0$, by $0$ for
$\beta < 0$), so $V^*_h(s)$ is finite and $V^\pi_h(s) \le V^*_h(s)$ for every policy $\pi$; hence
$Z^\pi_h(s) \le Z^*_h(s)$ ($\beta > 0$), resp.\ $Z^*_h(s) \le Z^\pi_h(s)$ ($\beta < 0$), for every
policy $\pi$. (The supremum is attained, by the greedy policy of Lemma~\ref{lem:opt_bellman}.)
\end{lemma}

\begin{proof}
\uses{def:opt_value, def:entropic_value, lem:exp_value_range}
\leanok
By Lemma~\ref{lem:exp_value_range}, $Z^\pi_h(s) > 0$ and $Z^\pi_h(s)$ lies between $1$ and
$e^{\beta(H - h)}$, so $V^\pi_h(s) = \beta^{-1} \log Z^\pi_h(s)$ is bounded above uniformly in
$\pi$ and the supremum $V^*_h(s)$ is a real number with $V^\pi_h(s) \le V^*_h(s)$. The map
$x \mapsto e^{\beta x}$ is increasing for $\beta > 0$ and decreasing for $\beta < 0$, and
$Z^\pi_h(s) = e^{\beta V^\pi_h(s)}$ (Definition~\ref{def:entropic_value}, using $Z^\pi_h(s) > 0$),
which gives the inequalities.
\end{proof}

\begin{lemma}[Optimal exponential Bellman equation]\label{lem:opt_bellman}
\lean{Learning.MDP.Episodic.optExpValueRec, Learning.MDP.Episodic.optExpQRec, Learning.MDP.Episodic.optPolicy, Learning.MDP.Episodic.expValue_optPolicy, Learning.MDP.Episodic.expValue_le_optExpValueRec, Learning.MDP.Episodic.optExpValueRec_le_expValue, Learning.MDP.Episodic.optEntropicValue_eq, Learning.MDP.Episodic.optExpValue_eq_optExpValueRec, Learning.MDP.Episodic.optExpValue_eq_expValue_optPolicy, Learning.MDP.Episodic.exists_expValue_eq_optExpValue, Learning.MDP.Episodic.exists_entropicValue_eq_optEntropicValue, Learning.MDP.Episodic.optExpValue_succ_eq, Learning.MDP.Episodic.optExpQ_le_optExpValue, Learning.MDP.Episodic.optExpValue_le_optExpQ, Learning.MDP.Episodic.optExpValue_of_le, Learning.MDP.Episodic.optExpQ_of_hasRewardFn, Learning.MDP.Episodic.optExpQ_of_hasRewardFn_eq_vecExp, Learning.MDP.Episodic.expQ_le_optExpQ, Learning.MDP.Episodic.optExpQ_le_expQ}\leanok
\uses{def:opt_value, lem:exp_bellman, lem:exp_value_monotone, lem:opt_exp_value_eq}
Let $\beta \ne 0$ and assume that the rewards of $\mathcal M$ are in $[0, 1]$ (for instance,
$\mathcal M$ has a reward function $r$ with values in $[0, 1]$), with finitely many actions and
countably many states. Then the supremum $V^*_h(s)$ is attained,
$Z^*_{h} = 1$ for $h \ge H + 1$, $U^*_h(s, a) = e^{\beta r_h(s, a)}(p_h Z^*_{h+1})(s, a)$, and for
every $h \le H$ and $s$,
\[
  Z^*_h(s) = \max_{a \in \As} U^*_h(s, a) \quad (\beta > 0),
  \qquad
  Z^*_h(s) = \min_{a \in \As} U^*_h(s, a) \quad (\beta < 0).
\]
Moreover the \emph{greedy policy} $\pi^*_h(s) \in \argmax_a U^*_h(s, a)$ ($\beta > 0$), resp.\
$\argmin_a U^*_h(s, a)$ ($\beta < 0$), satisfies $Z^{\pi^*}_h(s) = Z^*_h(s)$ for all $h$ and $s$:
it is optimal from every state at every step, and in particular $V^*_h(s) = V^{\pi^*}_h(s)$.
\end{lemma}

\begin{proof}
\uses{lem:exp_bellman, lem:exp_value_monotone, lem:opt_exp_value_eq, def:opt_value, lem:exp_value_range}
\leanok
Let $\beta > 0$ (the case $\beta < 0$ is identical with $\min$ in place of $\max$, $\argmin$ in
place of $\argmax$ and all inequalities reversed). Define the backward recursion
$\tilde Z_h(s) := 1$ for $h \ge H + 1$ and
$\tilde Z_h(s) := \max_a \tilde U_h(s, a)$ with
$\tilde U_h(s, a) := e^{\beta r_h(s, a)} (p_h \tilde Z_{h+1})(s, a)$ for $h \le H$ (the maximum
over the finite set of actions exists), and let $\pi^*_h(s) \in \argmax_a \tilde U_h(s, a)$ be a
greedy policy of this recursion.
\begin{enumerate}
  \item[(a)] $Z^{\pi^*} = \tilde Z$: by backward induction on $h$, using the exponential
        Bellman equation (Lemma~\ref{lem:exp_bellman}) for $\pi^*$ and the definition of $\pi^*_h$.
  \item[(b)] For every policy $\pi$, $Z^\pi_h(s) \le \tilde Z_h(s)$: by backward induction on $h$,
        $Z^\pi_h(s) = e^{\beta r_h(s, \pi_h(s))}(p_h Z^\pi_{h+1})(s, \pi_h(s))
        \le e^{\beta r_h(s, \pi_h(s))}(p_h \tilde Z_{h+1})(s, \pi_h(s)) = \tilde U_h(s, \pi_h(s))
        \le \tilde Z_h(s)$, by the Bellman equation, the induction hypothesis and monotonicity
        (Lemma~\ref{lem:exp_value_monotone}).
\end{enumerate}
By (b), $V^\pi_h(s) \le \beta^{-1}\log \tilde Z_h(s) = V^{\pi^*}_h(s)$ for every $\pi$, so the supremum
$V^*_h(s)$ is at most $V^{\pi^*}_h(s)$, and $V^{\pi^*}_h(s) \le V^*_h(s)$ by
Lemma~\ref{lem:opt_exp_value_eq}: the supremum is attained at $\pi^*$ and
$Z^* = Z^{\pi^*} = \tilde Z$. The Bellman equation and the formula for $U^*_h$ then follow from
the recursion defining $\tilde Z$ and $\tilde U$.
\end{proof}

\textbf{Lean remark.} The recursion is \texttt{optExpValueRec} / \texttt{optExpQRec} (a
well-founded recursion on \texttt{H - h}) and the greedy policy is
\texttt{optPolicy : ℕ → S → A}, built with \texttt{Finite.exists\_max} /
\texttt{Finite.exists\_min} on the finite action type and measurable on a countable discrete
state space (\texttt{measurable\_optPolicy}); (a) is \texttt{expValue\_optPolicy}, (b) is
\texttt{expValue\_le\_optExpValueRec} (\texttt{optExpValueRec\_le\_expValue} for $\beta < 0$),
and the conclusions are \texttt{optEntropicValue\_eq}, \texttt{optExpValue\_eq\_optExpValueRec}
and \texttt{optExpValue\_succ\_eq}. This is where the countability of the state space and the
finiteness of the action space enter (\texttt{[Countable S] [Finite A] [Nonempty A]}); the
supremum over all measurable policies needs no finiteness of the set of policies.

\begin{lemma}[Ranges of the optimal exponential values]\label{lem:opt_exp_value_range}
\lean{Learning.MDP.Episodic.optExpValue_mem_Icc, Learning.MDP.Episodic.optExpQ_mem_Icc, Learning.MDP.Episodic.optExpValue_pos}\leanok
\uses{def:opt_value, lem:exp_value_range, lem:opt_exp_value_eq}
For $\beta \ne 0$ and rewards in $[0, 1]$, for every $h \le H + 1$, $s$ and $a$,
\[
  Z^*_h(s) \in \big[\min(1, e^{\beta(H + 1 - h)}),\ \max(1, e^{\beta(H + 1 - h)})\big],
  \qquad
  U^*_h(s, a) \in \big[\min(1, e^{\beta(H + 1 - h)}),\ \max(1, e^{\beta(H + 1 - h)})\big]
\]
($h \le H$ for the second); in particular $Z^*_{h+1}$ takes values in $[1, e^{\beta(H - h)}]$
for $\beta > 0$ and in $[e^{\beta(H - h)}, 1]$ for $\beta < 0$, an interval of length
$\abs{e^{\beta(H - h)} - 1}$, and $Z^*_h > 0$, $U^*_h > 0$.
\end{lemma}

\begin{proof}
\uses{lem:exp_value_range, lem:opt_exp_value_eq, def:opt_value}
\leanok
By Lemma~\ref{lem:opt_exp_value_eq}, $Z^*_h(s) = Z^{\pi^\circ}_h(s)$ for some policy $\pi^\circ$, and
Lemma~\ref{lem:exp_value_range} gives the range. For $U^*_h(s, a)$, argue exactly as in the proof
of Lemma~\ref{lem:exp_value_range}: $e^{\beta R}$ lies between $1$ and $e^{\beta}$, $Z^*_{h+1}(S')$
between $1$ and $e^{\beta(H - h)}$, and the expectation of the product lies between $1$ and
$e^{\beta(H + 1 - h)}$.
\end{proof}

\begin{lemma}[Value gap in the exponential space]\label{lem:value_gap_log}
\lean{Learning.MDP.Episodic.optEntropicValue_sub_entropicValue, Learning.MDP.Episodic.optEntropicValue_sub_entropicValue_le_of_pos, Learning.MDP.Episodic.optEntropicValue_sub_entropicValue_le_of_neg, Learning.MDP.Episodic.optEntropicValue_eq_inv_mul_log}\leanok
\uses{def:opt_value, def:entropic_value, lem:exp_value_range, lem:opt_exp_value_eq}
Assume rewards in $[0, 1]$ and $\beta \ne 0$. For every policy $\pi$ and state $s$,
\[
  V^*_1(s) - V^\pi_1(s) = \frac1\beta \log \frac{Z^*_1(s)}{Z^\pi_1(s)}
  = \frac1\beta \log\Big(1 + \frac{Z^*_1(s) - Z^\pi_1(s)}{Z^\pi_1(s)}\Big) ,
\]
and consequently, for $\eps \ge 0$:
\begin{enumerate}
  \item[(i)] if $\beta > 0$ and $Z^*_1(s) - Z^\pi_1(s) \le (e^{\beta\eps} - 1)\, Z^\pi_1(s)$, then
        $V^*_1(s) - V^\pi_1(s) \le \eps$;
  \item[(ii)] if $\beta < 0$ and $Z^\pi_1(s) - Z^*_1(s) \le (1 - e^{\beta\eps})\, Z^\pi_1(s)$, then
        $V^*_1(s) - V^\pi_1(s) \le \eps$.
\end{enumerate}
The same holds at every step $h$ in place of $1$.
\end{lemma}

\begin{proof}
\uses{def:entropic_value, def:opt_value, lem:exp_value_range, lem:opt_exp_value_eq}
\leanok
$Z^\pi_1(s) > 0$ and $Z^*_1(s) > 0$ (Lemma~\ref{lem:exp_value_range},
Lemma~\ref{lem:opt_exp_value_range} through Lemma~\ref{lem:opt_exp_value_eq}), so
$V^*_1 - V^\pi_1 = \beta^{-1}(\log Z^*_1 - \log Z^\pi_1) = \beta^{-1}\log(Z^*_1/Z^\pi_1)$ and
$Z^*_1/Z^\pi_1 = 1 + (Z^*_1 - Z^\pi_1)/Z^\pi_1$. (i) The hypothesis is $Z^*_1(s) \le e^{\beta\eps} Z^\pi_1(s)$,
so $\log(Z^*_1/Z^\pi_1) \le \beta\eps$ and dividing by $\beta > 0$ gives the claim. (ii) The
hypothesis is $e^{\beta\eps} Z^\pi_1(s) \le Z^*_1(s)$, i.e.\ $Z^\pi_1/Z^*_1 \le e^{-\beta\eps} = e^{\abs\beta\eps}$,
so $V^*_1 - V^\pi_1 = \abs\beta^{-1}\log(Z^\pi_1/Z^*_1) \le \eps$ (for $\beta < 0$,
$\beta^{-1}\log(Z^*/Z^\pi) = \abs\beta^{-1}\log(Z^\pi/Z^*)$).
\end{proof}

\textbf{Lean remark.} This is the identity behind the stopping rule of Section~4 of the
paper (``Stopping rule'') and of its Appendix~B.2; both directions of the sign of $\beta$ are
needed. The formal statements hold at every step \texttt{h : ℕ} and only assume
$Z^\pi_h(s) > 0$ (not the range of the rewards); $Z^*_h(s) = e^{\beta V^*_h(s)} > 0$ holds by
definition.

\section{The maximal return}
\label{sec:mdp_gmax}

\begin{definition}[Maximal achievable return]\label{def:gmax}
\lean{Learning.MDP.Episodic.maxReturn}\leanok
\uses{def:step_law, def:episode_return}
For an initial state $s_1$, the \emph{maximal achievable return} of $\mathcal M$ (Definition~1 of
the paper) is
\[
  \Gmax(\mathcal M) := \sup_{\pi}\ \operatorname{ess\,sup}_{\Prob^\pi_{(s_1, 1)}} R^\pi_1 ,
\]
the supremum over the policies of the essential supremum of the return of an episode under the
trajectory law of the policy from $(s_1, 1)$.
\end{definition}

\textbf{Lean remark.} \texttt{maxReturn M H s₁ := ⨆ π : \{π : ℕ → S → A // ∀ k, Measurable (π k)\}, essSup (episodeReturn H) (stepLaw M π (s₁, 0))}
(a supremum over the measurable policies, bounded by $H$ for rewards in $[0, 1]$:
\texttt{bddAbove\_essSup\_episodeReturn}; \texttt{essSup} in $\R$, a conditionally complete
lattice, is the Mathlib \texttt{MeasureTheory.essSup}). The paper's ``supremum over all sources of
randomness'' is the essential supremum under the trajectory law: the two agree for the
finitely supported laws of the paper (finite state space, deterministic rewards) and the
essential supremum is the version that makes sense for reward kernels.

\begin{lemma}[The maximal return is between $0$ and $H$]\label{lem:gmax_le_horizon}
\lean{Learning.MDP.Episodic.maxReturn_nonneg, Learning.MDP.Episodic.maxReturn_le, Learning.MDP.Episodic.ae_episodeReturn_le_maxReturn, Learning.MDP.Episodic.essSup_episodeReturn_mem_Icc}\leanok
\uses{def:gmax, lem:exp_value_range}
For rewards in $[0, 1]$: $0 \le \Gmax(\mathcal M) \le H$, and for every policy $\pi$,
$R^\pi_1 \le \Gmax(\mathcal M)$ $\Prob^\pi_{(s_1, 1)}$-almost surely.
\end{lemma}

\begin{proof}
\uses{def:gmax, lem:exp_value_range}
\leanok
By Lemma~\ref{lem:exp_value_range}, $R^\pi_1 \in [0, H]$ a.s.\ under every $\Prob^\pi_{(s_1, 1)}$, so
each essential supremum lies in $[0, H]$ (the essential supremum of a function bounded a.s.\ by
$H$ is at most $H$, and it is at least $0$ since $R^\pi_1 \ge 0$ a.s.\ and the measure is a
probability measure); the supremum over the finite set of policies is then in $[0, H]$. The
last claim is the defining property of the essential supremum
($R^\pi_1 \le \operatorname{ess\,sup} R^\pi_1$ a.s., Mathlib \texttt{ae\_le\_essSup}) together with
$\operatorname{ess\,sup} R^\pi_1 \le \Gmax(\mathcal M)$ (\texttt{le\_ciSup}).
\end{proof}

\section{Episodes as rounds, BPI algorithms and the PAC property}
\label{sec:mdp_bpi}

The learner interacts with $\mathcal M$ in the online model: at each episode it chooses a
policy, plays it from the fixed initial state $s_1$ and observes the states of the episode.
Since the rewards are known and deterministic (Definition~\ref{def:reward_fn}), the state
sequence determines the episode; the whole interaction is an LML algorithm-environment
sequence in which one round is one episode.

\begin{definition}[Episode environment]\label{def:episode_env}
\lean{Learning.MDP.Episodic.statesLaw, Learning.MDP.Episodic.statesKernel, Learning.MDP.Episodic.statesEnv}\leanok
\uses{def:step_law, def:episode_return, def:policy}
Fix an initial state $s_1$. The \emph{episode environment} of $\mathcal M$ from $s_1$ is the
stationary LML environment (\texttt{stationaryEnv}) without observations, whose actions are the
policies $\pi \in \As^{\{1, \dots, H\} \times \Ss}$ and whose feedback to the action $\pi$ is the
state sequence $(S_1, \dots, S_{H+1})$ of an episode of $\pi$ from $s_1$, drawn from the
\emph{states kernel} $\nu_{\mathcal M}(\pi) :=$ the law of $(S_1, \dots, S_{H+1})$ under
$\Prob^\pi_{(s_1, 1)}$. A run of an algorithm in this environment is a sequence
$(X_t, Y_t)_{t \ge 0}$ where $X_t = \pi^{t+1}$ is the policy of the episode $t + 1$ and
$Y_t = (s^{t+1}_1, \dots, s^{t+1}_{H+1})$ its state sequence; we write $a^{t+1}_h := \pi^{t+1}_h(s^{t+1}_h)$
for the action played at step $h$ of the episode $t + 1$.
\end{definition}

\textbf{Lean remark.} \texttt{statesLaw M H s₁ π := (stepLaw M π (s₁, 0)).map (episodeStates H)}
for a policy \texttt{π : ℕ → S → A}, \texttt{statesKernel M H s₁ : Kernel (Policy S A H) (Traj S H)},
\texttt{π ↦ statesLaw M H s₁ π.extend} (measurable on the countable discrete type of
$H$-step policies), and \texttt{statesEnv M H s₁ := stationaryEnv (statesKernel M H s₁) : Environment Unit (Policy S A H) (Traj S H)}.
The observation type is \texttt{Unit}. The rewards are not part of the feedback: an
algorithm for this environment is a learner who knows $r$; this is the reason why the PAC
property below is stated for the class $\mathcal M_r$.

\begin{lemma}[Law of an episode given the past]\label{lem:episode_env_law}
\lean{Learning.MDP.Episodic.hasCondDistrib_feedback_history_action_statesEnv, Learning.MDP.Episodic.hasCondDistrib_feedback_statesEnv}\leanok
\uses{def:episode_env, lem:step_law_markov}
Let $(X, Y)$ be an algorithm-environment sequence for any algorithm and the episode
environment of $\mathcal M$ from $s_1$, on $(\Omega, \Prob)$. For every $t$, the conditional law of
$Y_t$ given the history $(X_i, Y_i)_{i < t}$ and $X_t$ is $\nu_{\mathcal M}(X_t)$: for every
bounded measurable $F$ on $\Ss^{H+1}$,
$\E[F(Y_t) \mid (X_i, Y_i)_{i < t}, X_t] = \E^{X_t}_{(s_1, 1)}[F(S_1, \dots, S_{H+1})]$ a.s. The
within-episode form (the law of $s^{t+1}_{h+1}$ given the past and $s^{t+1}_{\le h}$) is
Lemma~\ref{lem:pmdp_episode_env_step}.
\end{lemma}

\begin{proof}
\uses{def:episode_env}
\leanok
This is the defining property of an algorithm-environment sequence for a stationary environment
(LML \texttt{IsObliviousEnv.hasCondDistrib\_feedback\_history\_action} together with
\texttt{feedbackCondAction\_stationaryEnv}: the feedback kernel of \texttt{stationaryEnv ν} at every
round is $\nu$ applied to the action).
\end{proof}

\textbf{Lean remark.} \texttt{hasCondDistrib\_feedback\_history\_action\_statesEnv}:
\texttt{HasCondDistrib (Y t) (fun ω ↦ ((history O X Y t ω, O t ω), X t ω)) ((statesKernel M s₁).prodMkLeft \_) P},
and \texttt{hasCondDistrib\_feedback\_statesEnv} given the policy only.

\begin{definition}[BPI algorithms, runs and stopping times]\label{def:bpi_alg}
\lean{Learning.MDP.Episodic.BPIAlg}\leanok
\uses{def:episode_env}
A \emph{best-policy identification (BPI) algorithm} is an LML identification algorithm
(\texttt{IdentAlg}) in the episode environment: a triple $\mathcal A = (\mathrm{alg}, \mathcal T, \rho)$
of a \emph{sampling rule} $\mathrm{alg}$, an LML algorithm with actions in the policies and
feedback in $\Ss^{H+1}$ (at the round $t$ it chooses the policy $\pi^{t+1}$ of the episode
$t + 1$ as a function, possibly randomized, of the history of the first $t$ episodes), a
\emph{stopping rule} $\mathcal T$, a measurable set of finite histories of variable length
(``stop after $t$ episodes when their history belongs to $\mathcal T$''), and an \emph{output
rule} $\rho$, a Markov kernel from histories of variable length to policies. A \emph{run} of
$\mathcal A$ in the episode environment of $\mathcal M$ from $s_1$ on $(\Omega, \Prob)$ is a triple
$(X, Y, \hat\pi)$ where $(X, Y)$ is an algorithm-environment sequence for $\mathrm{alg}$ and the
episode environment and the output $\hat\pi$ has conditional law $\rho(\hist_\tau)$ given the
history $\hist_\tau$ at the stopping time. The \emph{number of episodes} (sample complexity) is
the stopping time
\[
  \tau := \min\{t \in \N : (X_i, Y_i)_{i < t} \in \mathcal T\} \in \N \cup \{\infty\} ,
\]
a stopping time of the history filtration (LML \texttt{isStoppingTime\_stoppingTime}); the law of
$(\hist_\tau, \hat\pi)$ is the same in all runs, and we write $\Prob_{\mathcal M}$ for it when a
statement only depends on it (LML \texttt{IdentAlg.outputMeasure} for the law of $\hat\pi$).
\end{definition}

\textbf{Lean remark.} \texttt{BPIAlg S A H := IdentAlg Unit (Policy S A H) (Traj S H) (Policy S A H)};
the stopping rule is the field \texttt{stopSet : Set (Σ n, Hist Unit (Policy S A H) (Traj S H) n)},
the output rule \texttt{output : Kernel (Σ n, Hist \_ \_ \_ n) (Policy S A H)}, the stopping time
\texttt{IdentAlg.stoppingTime A O X Y : Ω → ℕ∞} (\texttt{hittingAfter} of the process of
histories \texttt{sigmaHistory}), runs are \texttt{IdentAlg.IsRun A env O X Y out P} (with
\texttt{O : ℕ → Ω → Unit} trivial), and \texttt{IdentAlg.stoppedHist} is $\hist_\tau$
(\texttt{stoppedHist\_mem\_stopSet\_of\_ne\_top}: it belongs to the stopping rule when
$\tau < \infty$). When $\tau = \infty$ the stopped history is a history of an arbitrary
number of rounds; all statements about the output are on the event $\{\tau < \infty\}$ or
about the law of the output.

\begin{definition}[$(\eps, \delta)$-PAC algorithms for entropic BPI]\label{def:entropic_pac}
\lean{Learning.MDP.Episodic.IsEntropicPAC}\leanok
\uses{def:bpi_alg, def:opt_value, def:entropic_value, def:reward_fn}
Let $r$ be a reward function, $\beta \ne 0$, $\eps > 0$, $\delta \in (0, 1)$ and $s_1 \in \Ss$. A
BPI algorithm $\mathcal A$ is \emph{$(\eps, \delta)$-PAC for entropic BPI with reward function
$r$, parameter $\beta$ and initial state $s_1$} (Definition~2 of the paper) if for every MDP
$\mathcal M \in \mathcal M_r$, the output $\hat\pi$ of $\mathcal A$ in the episode environment of
$\mathcal M$ from $s_1$ satisfies
\[
  \Prob_{\mathcal M}\big( V^*_1(s_1) - V^{\hat\pi}_1(s_1) > \eps \big) \le \delta ,
\]
where the values are those of $\mathcal M$. Equivalently: in every run of $\mathcal A$ in the episode
environment of every $\mathcal M \in \mathcal M_r$, $\Prob(V^*_1(s_1) - V^{\hat\pi}_1(s_1) \le \eps) \ge 1 - \delta$.
\end{definition}

\textbf{Lean remark.} \texttt{IsEntropicPAC 𝒜 r β ε δ s₁ := 𝒜.IsPAC (fun M : \{M // M.HasRewardFn r\} ↦ statesEnv M.1 s₁) (fun M π ↦ ε < optEntropicValue M.1 β 0 s₁ - entropicValue M.1 β π 0 s₁) δ}.
LML's \texttt{IdentAlg.IsPAC env bad δ} says that, for every parameter $\theta$ (here every
$\mathcal M \in \mathcal M_r$), the law \texttt{outputMeasure} of the output in \texttt{env θ} gives
mass at most $\delta$ to the bad outputs; for a run $(X, Y, \hat\pi)$ on $(\Omega, \Prob)$,
\texttt{IsPAC.measureReal\_bad\_of\_isRun} gives $\Prob(\hat\pi \text{ bad}) \le \delta$
(\texttt{IsRun.hasLaw\_output}: the output of a run has law \texttt{outputMeasure}). Conversely,
to prove that an algorithm is PAC one proves the bound for one run, the canonical run on the
Ionescu--Tulcea space (LML \texttt{IdentAlg.runMeasure}, \texttt{isRun\_runMeasure}), which
is what Theorem~\ref{thm:upper_bound_pac} does. Two points of the encoding are decisions
recorded in the overview (deviation~7): the PAC property is a property of the \emph{algorithm}
(quantified over the environments of the class), not of a run, and the class is
$\mathcal M_r$ for the \emph{known} reward function $r$ of the algorithm and arbitrary
transition kernels; the lower bound (Theorem~\ref{thm:lower_bound}) only assumes PAC for the
class with the reward function of its hard instances, a weaker hypothesis than PAC for all
reward functions, hence a stronger theorem. Note also that the paper's Definition~2 does not
require $\tau < \infty$: an algorithm that never stops has an output law given by the output
kernel applied to an infinite history in LML, and finiteness of $\tau$ is a separate
conclusion (Theorem~\ref{thm:upper_bound_complexity}).
